% ------------------------------------------------------------------------------------
%                           ESERCIZI TRASMISSIONE e MODULAZIONE
% ------------------------------------------------------------------------------------

\chapter{Esercizi in Aula - Trasmissione e Modulazione}

\section{Modulazione PWM}

Vediamo un esempio di modulazione di
tipo \textbf{PWM}: \textbf{Pulse Width Modulation}.

\begin{remark}[PWM]
    In questo tipo di modulazione si utilizzano \textbf{onde quadre} con
    \textbf{diverse larghezze degli impulsi}.

    E' ampiementa utilizzata in \textbf{elettronica digitale} perchè
    permette di \textit{controllare la potenza di motori elettrici}
    usando \textit{segnali digitali}.
\end{remark}

Usiamo per esempio $M=4$ con i segnali seguenti, mostrati 
anche in figura~\ref{fig:B_01_EsempioPWM}.
\[
    \begin{aligned}
        s_{1}(t) &= \begin{cases}
            A & 0 < t <\frac{T}{4} \\ 
            0 & \text{altrimenti}
        \end{cases} \\ 
        s_{2}(t) &= \begin{cases}
            A & 0 < t <\frac{T}{2} \\ 
            0 & \text{altrimenti}
        \end{cases} \\ 
        s_{3}(t) &= \begin{cases}
            A & 0 < t <\frac{3T}{4} \\ 
            0 & \text{altrimenti}
        \end{cases} \\ 
        s_{4}(t) &= \begin{cases}
            A & 0 < t < T \\ 
            0 & \text{altrimenti}
        \end{cases} \\ 
    \end{aligned}
\]

\begin{figure}[!ht]
    \centering
    \pgfplotsset{
        pwmStyle/.style={
            width=6.5cm,
            height=3.5cm,
            axis lines=middle,
            xmin=-0.05, xmax=1.15,
            ymin=-0.1, ymax=1.3,
            xtick={0, 0.25, 0.5, 0.75, 1},
            xticklabels={\(0\), \(\frac{T}{4}\), \(\frac{T}{2}\), \(\frac{3T}{4}\), \(T\)},
            ytick={0, 1},
            yticklabels={\(0\), \(A\)},
            ylabel style={at={(0,1)}, anchor=south east},
            grid=major,
            grid style={dashed, gray!20},
            axis line style={-Stealth, thick},
        }
    }

    \begin{tikzpicture}

        % --- s1(t) ---
        \begin{axis}[name=ax1, pwmStyle, ylabel={\(s_1(t)\)}]
            \addplot[very thick, nord10] coordinates {(0,0) (0,1) (0.25,1) (0.25,0) (1,0)};
        \end{axis}

        % --- s2(t) ---
        \begin{axis}[name=ax2, pwmStyle, at={(ax1.east)}, anchor=west, xshift=1cm, ylabel={\(s_2(t)\)}]
            \addplot[very thick, nord10] coordinates {(0,0) (0,1) (0.5,1) (0.5,0) (1,0)};
        \end{axis}

        % --- s3(t) ---
        \begin{axis}[name=ax3, pwmStyle, at={(ax1.south west)}, 
            anchor=north west, yshift=-1cm, ylabel={\(s_3(t)\)}]
            \addplot[very thick, nord10] coordinates {(0,0) (0,1) (0.75,1) (0.75,0) (1,0)};
        \end{axis}

        % --- s4(t) ---
        \begin{axis}[name=ax4, pwmStyle, at={(ax2.south west)}, 
            anchor=north west, yshift=-1cm, ylabel={\(s_4(t)\)}]
            \addplot[very thick, nord10] coordinates {(0,0) (0,1) (1,1) (1,0)};
        \end{axis}
    \end{tikzpicture}
    \caption{Esempio di forme d'onda per modulazione PWM}
    \label{fig:B_01_EsempioPWM}
\end{figure}

Supponiamo di voler \textit{ricavare una base ortonormale} per questi segnali.
Notiamo subito che:

\begin{remark}
    I segnali sono tutti \textbf{linearmente indipendenti} (non siamo 
    in grado di scrivere uno come combinazione lineare degli altri), 
    tuttavia \textbf{non sono ortogonali}.
\end{remark}

Per ottenere una \textbf{base ortonormale} potremmo usare il
procedimento di Gram-Schmidt, tuttavia risulta spesso particolarmente
laborioso. Possiamo invece \textbf{scegliere intelligentemente} una base
del tipo:
\[
    \begin{aligned}
        \phi_{1}(t) &= \begin{cases}
            A' & 0 < t <\frac{T}{4} \\ 
            0 & \text{altrimenti}
        \end{cases} \\ 
        \phi_{2}(t) &= \begin{cases}
            A' & \frac{T}{4} < t <\frac{T}{2} \\ 
            0 & \text{altrimenti}
        \end{cases} \\ 
        \phi_{3}(t) &= \begin{cases}
            A' & \frac{T}{2} < t <\frac{3T}{4} \\ 
            0 & \text{altrimenti}
        \end{cases} \\ 
        \phi_{4}(t) &= \begin{cases}
            A' & \frac{3T}{4} < t < T \\ 
            0 & \text{altrimenti}
        \end{cases} \\ 
    \end{aligned}
\]

I segnali di questa base sono mostrati in figura~\ref{fig:B_02_EsempioPWMBaseOrto}.

\begin{figure}[!ht]
    \centering
    \pgfplotsset{
        basisStyle/.style={
            width=6.5cm,
            height=3.5cm,
            axis lines=middle,
            xmin=-0.05, xmax=1.05,
            ymin=-0.1, ymax=1.3,
            xtick={0, 0.25, 0.5, 0.75, 1},
            xticklabels={\(0\), \(\frac{T}{4}\), \(\frac{T}{2}\), \(\frac{3T}{4}\), \(T\)},
            ytick={0, 1},
            yticklabels={\(0\), \(A'\)},
            ylabel style={at={(0,1)}, anchor=south east},
            grid=major,
            grid style={dashed, gray!20},
            axis line style={-Stealth, thick},
        }
    }

    \begin{tikzpicture}

        % --- phi_1(t) ---
        \begin{axis}[name=ax5, basisStyle, ylabel={\(\phi_1(t)\)}]
            \addplot[very thick, nord14] coordinates {(0,0) (0,1) (0.25,1) (0.25,0) (1,0)};
        \end{axis}

        % --- phi_2(t) ---
        \begin{axis}[name=ax6, basisStyle, at={(ax5.east)}, anchor=west, xshift=1cm, ylabel={\(\phi_2(t)\)}]
            \addplot[very thick, nord14] coordinates {(0,0) (0.25,0) (0.25,1) (0.5,1) (0.5,0) (1,0)};
        \end{axis}

        % --- phi_3(t) ---
        \begin{axis}[name=ax7, basisStyle, at={(ax5.south west)}, 
            anchor=north west, yshift=-1cm, ylabel={\(\phi_3(t)\)}]
            \addplot[very thick, nord14] coordinates {(0,0) (0.5,0) (0.5,1) (0.75,1) (0.75,0) (1,0)};
        \end{axis}

        % --- phi_4(t) ---
        \begin{axis}[name=ax8, basisStyle, at={(ax6.south west)}, 
            anchor=north west, yshift=-1cm, ylabel={\(\phi_4(t)\)}]
            \addplot[very thick, nord14] coordinates {(0,0) (0.75,0) (0.75,1) (1,1) (1,0)};
        \end{axis}

    \end{tikzpicture}
    \caption{Base scelta per la PWM d'esempio}
    \label{fig:B_02_EsempioPWMBaseOrto}
\end{figure}

Infine notiamo che i segnali sono tutti a \textbf{supporto disgiunto}, per cui il loro
\textbf{prodotto interno} è \textbf{sicuramente nullo} e quindi \textbf{formano una base ortogonale}. 
La scelta è stata intelligente perchè abbiamo banalmente:
\[
    \begin{aligned}
        s_{1}(t) &= \alpha \phi_{1}(t) \\ 
        s_{2}(t) &= \alpha \left[ \phi_{1}(t) + \phi_{2}(t) \right] \\ 
        s_{3}(t) &= \alpha \left[ \phi_{1}(t) + \phi_{2}(t) + \phi_{3}(t) \right] \\ 
        s_{4}(t) &= \alpha \left[ \phi_{1}(t) + \phi_{2}(t) + \phi_{3}(t) + \phi_{4}(t) \right] \\ 
    \end{aligned}
\]

Ci resta infine solo da \textbf{determinare l'ampiezza} affinchè la norma dei vettori 
sia unitaria. Per tutti e quattro i segnali l'energia risulta:
\[
    E_{\phi} = A'^2 \frac{T}{4}
\]

Per cui scegliamo:
\[
    A' = \frac{2}{\sqrt{T}} = \frac{2\sqrt{T}}{T}
\]

\section{Canale con attenuazione e ritardo}

Consideriamo ora un esempio di canale che introduce sia \textbf{un'attenuazione} che
\textbf{un ritardo} nel \textbf{segnale trasmesso}.
Possiamo supporre che la trasformazione introdotta dal canale sia:
\[
    g_c(t) = \delta(t) + \frac{1}{2}\delta \left( t - \frac{T}{4} \right)
\]

\begin{remark}[ATTENUAZIONE E RITARDO]
    Notiamo che la \textit{risposta} del canale è data da \textit{due termini}:
    \begin{itemize}
        \item Il primo ($\delta(t)$) corrisponde al \textbf{cammino diretto}: il segnale 
            \textit{arriva diretto}, \textit{senza ritardo} e \textit{senza attenuazione}.
        \item Il secondo ($0.5\delta(t - T/4)$) corrisponde ad un 
            \textbf{cammino più lungo}: durante la trasmissione il segnale 
            \textit{colpisce ostacoli} e \textit{viene riflesso più volte}.
            Di conseguenza, \textbf{perde potenza} (fattore $1/2$) e arriva con un
            certo \textbf{ritardo} ($t-T/4$).
    \end{itemize}
\end{remark}

\begin{figure}[!ht]
    \centering
    \begin{tikzpicture}[
        >=Stealth,
        thick,
        % Stile per le antenne/trasmettitori
        antenna/.style={
            draw=nord0, 
            fill=nord14!30, 
            rectangle, 
            minimum width=1.5cm, 
            minimum height=1cm, 
            align=center, 
            font=\bfseries,
            text=nord0
        },
        % Stile per l'ostacolo
        obstacle/.style={
            draw=nord3, 
            fill=nord3!20, 
            rectangle, 
            minimum width=2.5cm, 
            minimum height=1.5cm, 
            align=center, 
            font=\small,
            text=nord3,
            rounded corners=2pt
        }
        ]

        % --- Nodi (Tx, Rx, Ostacolo) ---
        \node[antenna] (tx) at (0,0) {Tx};
        \node[antenna] (rx) at (8,0) {Rx};
        \node[obstacle] (obs) at (4, 3) {Ostacolo\\(Riflessione)};

        % --- Cammino Diretto (LOS) ---
        \draw[->, very thick, nord10] (tx.east) -- (rx.west) 
            node[midway, below, font=\small] {Cammino Diretto};
        % Etichetta matematica per il cammino diretto
        \node[below=0.6cm of tx.east, xshift=1.5cm, font=\scriptsize, text=nord10] {Ampiezza: \(1\)};
        \node[below=0.6cm of tx.east, xshift=3.5cm, font=\scriptsize, text=nord10] {Ritardo: \(0\)};

        % --- Cammino Riflesso ---
        \draw[->, thick, nord11, dashed] (tx.north) -- (obs.south west) -- (rx.north) 
            node[pos=0.45, below, sloped, font=\small] {Cammino Riflesso};

        % Etichette matematiche per il cammino riflesso
        \node[right=0.1cm of obs.south east, font=\scriptsize, text=nord11, align=left] {
                Attenuazione: \(\frac{1}{2}\)\\
                Ritardo: \(\frac{T}{4}\)
            };
    \end{tikzpicture}
    \caption{Rappresentazione del fenomeno di multicammino}
    \label{fig:B_03_MultipathChannel}
\end{figure}

Immaginiamo di trasmettere il segnale usando una \textbf{modulazione} che porta ad 
una \textbf{base ortonormale} $\mathcal{B} = \{\phi_{1}(t), \phi_{2}(t) , \phi_3(t), \phi_{4}(t)\}$, 
come quella definita sotto:
\[
    \begin{aligned}
        \phi_{1}(t) &= \begin{cases}
            \frac{2}{\sqrt{T}} & 0 < t <\frac{T}{4} \\ 
            0 & \text{altrimenti}
        \end{cases} \\ 
        \phi_{2}(t) &= \begin{cases}
            \frac{2}{\sqrt{T}} & 0 < t <\frac{T}{2} \\ 
            0 & \text{altrimenti}
        \end{cases} \\ 
        \phi_{3}(t) &= \begin{cases}
            \frac{2}{\sqrt{T}} & 0 < t <\frac{3T}{4} \\ 
            0 & \text{altrimenti}
        \end{cases} \\ 
        \phi_{4}(t) &= \begin{cases}
            \frac{2}{\sqrt{T}} & 0 < t < T \\ 
            0 & \text{altrimenti}
        \end{cases} \\ 
    \end{aligned}
\]

Notiamo che è la stessa base che avevamo usato nell'esempio precedente, ed è
mostrata in figura~\ref{fig:B_02_EsempioPWMBaseOrto}.

Consideriamo ora la \textbf{trasmissione} di un \textbf{singolo simbolo} $s(t)$ 
usando la base data:
\[
    s(t) = \sum_{i=1}^4 s_i \phi_i(t)
\]

Dove $s_{1},s_{2},s_{3},s_{4}$ sono le \textit{coordinate} di $s(t)$ rispetto alla base.

Il \textbf{simbolo ricevuto} $y(t)$ sarà la \textit{convoluzione} tra il segnale e la
\textit{risposta impulsiva del canale}:
\[
    y(t) = s(t) \ast g_c(t) = s(t) \ast 
    \left[ \delta(t) + \frac{1}{2}\delta \left( t - \frac{T}{4} \right) \right]
\]
Per la \textit{proprietà rivelatrice} questo diventa semplicemente:
\[
    y(t) = s(t) + \frac{1}{2}s \left( t - \frac{T}{4} \right)
\]

\begin{remark}
    Come ci aspettavamo alla luce dell'osservazione fatta sopra, riceviamo
    il \textbf{segnale originale} ed una sua \textbf{copia ritardata e attenuata}.
\end{remark}

Per trovare la \textbf{rappresentazione vettoriale} del segnale ricevuto $r(t)$
dobbiamo \textit{proiettare} $y(t)$ su \textit{ciascun elemento della base} $\phi_j(t)$:
\[
    \begin{aligned}
        r_j &= \langle y(t), \phi_j(t) \rangle = 
              \left\langle s(t) + \frac{1}{2}s\left( t - \frac{T}{4} \right), \phi_j (t) \right\rangle \\ 
            &= \underbrace{ \langle s(t), \phi_j(t) \rangle }_{\text{componente diretta}}
            + \frac{1}{2}\underbrace{\left\langle s \left( t - \frac{T}{4} \right), \phi_j(t) \right\rangle}_{
            \text{componente riflessa}}
    \end{aligned}
\]

Notiamo a questo punto che:
\begin{itemize}
    \item Il \textit{primo termine} è semplicemente $s_j$ (per l'ortogonalità della base)
    \item Per il \textit{secondo termine}, dobbiamo capire su quale $\phi_j$ \textit{"cade"}
        il \textbf{segnale ritardato}:
        \begin{itemize}
            \item $\phi_{1}(t)$ è \textit{"attivo"} per $t\in [0, T/4]$. Se lo ritardiamo
                di $T/4$ diventa \textit{"attivo"} per $t\in[T/4, T/2]$ e
                \textbf{coincide} con $\phi_{2}(t)$.
            \item Similmente, $\phi_{2}(t)$ ritardato cade su $\phi_{3}(t)$.
            \item E $\phi_{3}(t)$ ritardato cade su $\phi_{4}(t)$.
            \item La questione su $\phi_{4}$ è \textit{più complessa} ed
                è descritta sotto.
        \end{itemize}
\end{itemize}

Tralasciando $\phi_{4}$, notiamo che ciascun \textbf{segnale ritardato} $\phi_i(t - T/4)$ è
\textbf{coincidente} con l'\textbf{elemento successivo della base} $\phi_{i+1}(t)$.

Possiamo quindi scrivere il \textbf{simbolo ricevuto} in \textbf{forma vettoriale}
$\mathbf{r} = [ r_{1}, r_{2}, r_{3}, r_{4} ]^\top $ in funzione 
del \textit{simbolo trasmesso} $s = [s_{1}, s_{2}, s_{3}, s_{4}]^\top$:
\[
    \begin{cases}
        r_{1} = 1\cdot s_{1} + 0\cdot s_{2} + 0\cdot s_{3} + 0\cdot s_{4} \\
        r_{2} = \frac{1}{2}\cdot s_{1} + 1\cdot s_{2} + 0\cdot s_{3} + 0\cdot s_{4} \\
        r_{3} = 0\cdot s_{1} + \frac{1}{2}\cdot s_{2} + 1\cdot s_{3} + 0\cdot s_{4} \\
        r_{1} = 0\cdot s_{1} + 0\cdot s_{2} + \frac{1}{2}\cdot s_{3} + 1\cdot s_{4} \\
    \end{cases}
\]

O, in \textit{forma matriciale} (più compatta):
\[
    \mathbf{r} = H \mathbf{s} = 
    \begin{bmatrix} 1 & 0 & 0 & 0 \\ 
    \frac{1}{2} & 1 & 0 & 0 \\ 
    0 & \frac{1}{2} & 1 & 0 \\ 
    0 & 0 & \frac{1}{2} & 1
    \end{bmatrix} \begin{bmatrix} s_{1} \\ s_{2} \\ s_{3} \\ s_{4} \end{bmatrix} 
\]

Detto ciò, dedichiamo una discussione a parte a ciò che succede con $\phi_{4}(t)$.
Se ritardiamo $\phi_{4}(t)$ di $T/4$, otteniamo un segnale definito per
$t\in [T, 5T/4]$: questo intervallo esce al di \textbf{fuori del periodo di simbolo corrente},
andando a \textbf{sforare nel successivo} (in particolare, si sovrappone alla componente
relativa a $\phi_{1}(t)$ del \textit{simbolo successivo}).

\begin{remark}[ISI]
    Questo fenomeno è il cuore di un \textbf{problema} molto importante
    nelle telecomunicazioni: la \textbf{Inter-Symbol Interference} (\textbf{ISI}).
\end{remark}

Ciò significa che, se volessimo descrivere \textbf{l'intero flusso di dati} in modo completo, 
dovremmo considerare che il segnale ricevuto al tempo $t \in [T,5T/4]$ non dipende solo 
dal primo simbolo del nuovo blocco, ma è la \textbf{somma di}:
\begin{itemize}
    \item La \textbf{prima componente} del \textbf{nuovo simbolo trasmesso} $s_{1}^{(k)}$.
    \item L'\textbf{"eco"} dell'\textbf{ultima componente} del 
        \textbf{simbolo precedente} $\frac{1}{2}s_{4}^{(k-1)}$.
\end{itemize}
